\chapter*{Appendix}
\addcontentsline{toc}{chapter}{Appendix}

The full summary statistics (mean, standard deviation, median, \acrshort{iqr}, minimum, maximum, quartiles, and sample counts) per scenario, algorithm, and metric 
are available in the supplementary materials at \path{Code/Results/summary_stats.csv}.

The full raw benchmark data can be found at \path{Code/Results/raw_runs.csv}, and the corresponding string pairs with their LCS values are stored in \path{Code/Results/case_lcs_values.csv}.

% Write your appendix here.

\section*{Slope Ratios}
\label{app:sloperatios}
The following tables present the slope ratios between the \acrshort{sam} and the \acrshort{esa} for each scenario and metric, 
along with the coefficient of determination ($R^2$) as a measure of goodness of fit for the linear regression.

\begin{itemize}
    \item \textbf{SAM slope} and \textbf{ESA slope}: Fitted slope of the linear regression of the median measurement against string length for the respective algorithm. 
    The slope is to be interpreted as the constant $a$ in the linear model $y = a \cdot x + b$, where $y$ is the median measurement (e.g., build time) and $x$ is the length of one of the strings in the string pair.
    \item \textbf{SAM / ESA ratio}: Ratio between the two slopes. Values below $1$ mean that the \acrshort{sam} grows more slowly than the \acrshort{esa}, while values above $1$ mean the opposite.
    \item \textbf{SAM $R^2$} and \textbf{ESA $R^2$}: A measure of how closely a linear model fits the observed median values for each algorithm, with values closer to $1$ indicating a stronger linear relationship.
    \item \textbf{Factor}: An expression of how many times larger the greater slope is compared to the smaller one, always reported as a value $\geq 1$. 
    For example, a factor of $17.30\times$ means one algorithm's growth rate is 17.30 times steeper than the other's.
    \item \textbf{Worse}: The algorithm with the steeper slope. As in, the one that scales less favourably with increasing input length.
\end{itemize}

% Sorted and formatted benchmark tables
\subsection*{Build Time}
\begin{benchmarktable}{Slope Ratios: Build Time}{Linear regression slope fitted to median build time vs. string length }
\begin{tabularx}{\textwidth}{|X|X|X|X|X|X|X|X|}
\toprule
Scenario & SAM slope & ESA slope & SAM / ESA ratio & SAM $R^2$ & ESA $R^2$ & Factor & Worse \\
\midrule
Disjoint Alphabet & 0.0048 & 0.0825 & 0.0578 & 0.9996 & 0.9997 & $17.30\times$ & ESA \\
\hline
Mutated Implant & 0.0058 & 0.0868 & 0.0665 & 0.9993 & 0.9997 & $15.05\times$ & ESA \\
\hline
Near Identical & 0.0042 & 0.0915 & 0.0462 & 0.9989 & 0.9994 & $21.63\times$ & ESA \\
\hline
Random Uniform & 0.0043 & 0.0832 & 0.0511 & 0.9990 & 0.9996 & $19.57\times$ & ESA \\
\hline
Repetitive With Noise & 0.0062 & 0.0878 & 0.0709 & 0.9992 & 0.9997 & $14.11\times$ & ESA \\
\bottomrule
\end{tabularx}
\end{benchmarktable}

\subsection*{Query Time}
\begin{benchmarktable}{Slope Ratios: Query Time}{Linear regression slope fitted to median query time vs. string length }
\begin{tabularx}{\textwidth}{|X|X|X|X|X|X|X|X|}
\toprule
Scenario & SAM slope & ESA slope & SAM / ESA ratio & SAM $R^2$ & ESA $R^2$ & Factor & Worse \\
\midrule
Disjoint Alphabet & 0.0009 & 0.0029 & 0.3233 & 0.9998 & 0.9999 & $3.09\times$ & ESA \\
\hline
Mutated Implant & 0.0013 & 0.0030 & 0.4270 & 0.9986 & 0.9999 & $2.34\times$ & ESA \\
\hline
Near Identical & 0.0077 & 0.0031 & 2.4885 & 0.9979 & 0.9999 & $2.49\times$ & SAM \\
\hline
Random Uniform & 0.0012 & 0.0030 & 0.4036 & 0.9995 & 0.9999 & $2.48\times$ & ESA \\
\hline
Repetitive With Noise & 0.0012 & 0.0030 & 0.3913 & 0.9981 & 1.0000 & $2.56\times$ & ESA \\
\bottomrule
\end{tabularx}
\end{benchmarktable}

\subsection*{Total Time}
\begin{benchmarktable}{Slope Ratios: Total Time}{Linear regression slope fitted to median total time vs. string length }
\begin{tabularx}{\textwidth}{|X|X|X|X|X|X|X|X|}
\toprule
Scenario & SAM slope & ESA slope & SAM / ESA ratio & SAM $R^2$ & ESA $R^2$ & Factor & Worse \\
\midrule
Disjoint Alphabet & 0.0057 & 0.0854 & 0.0668 & 0.9997 & 0.9997 & $14.97\times$ & ESA \\
\hline
Mutated Implant & 0.0070 & 0.0898 & 0.0783 & 0.9994 & 0.9997 & $12.77\times$ & ESA \\
\hline
Near Identical & 0.0119 & 0.0946 & 0.1256 & 0.9985 & 0.9995 & $7.96\times$ & ESA \\
\hline
Random Uniform & 0.0054 & 0.0862 & 0.0632 & 0.9993 & 0.9996 & $15.83\times$ & ESA \\
\hline
Repetitive With Noise & 0.0074 & 0.0907 & 0.0814 & 0.9992 & 0.9997 & $12.29\times$ & ESA \\
\bottomrule
\end{tabularx}
\end{benchmarktable}

\subsection*{Index Size}
\begin{benchmarktable}{Slope Ratios: Index Size}{Linear regression slope fitted to median index size vs. string length }
\begin{tabularx}{\textwidth}{|X|X|X|X|X|X|X|X|}
\toprule
Scenario & SAM slope & ESA slope & SAM / ESA ratio & SAM $R^2$ & ESA $R^2$ & Factor & Worse \\
\midrule
Disjoint Alphabet & 0.3550 & 0.0892 & 3.9794 & 0.9999 & 1.0000 & $3.98\times$ & SAM \\
\hline
Mutated Implant & 0.4031 & 0.0892 & 4.5193 & 1.0000 & 1.0000 & $4.52\times$ & SAM \\
\hline
Near Identical & 0.3347 & 0.1159 & 2.8887 & 0.9998 & 1.0000 & $2.89\times$ & SAM \\
\hline
Random Uniform & 0.3347 & 0.0892 & 3.7522 & 0.9998 & 1.0000 & $3.75\times$ & SAM \\
\hline
Repetitive With Noise & 0.4430 & 0.0892 & 4.9661 & 1.0000 & 1.0000 & $4.97\times$ & SAM \\
\bottomrule
\end{tabularx}
\end{benchmarktable}

\subsection*{Peak Memory Usage}
\begin{benchmarktable}{Slope Ratios: Peak Memory}{Linear regression slope fitted to median peak memory vs. string length }
\begin{tabularx}{\textwidth}{|X|X|X|X|X|X|X|X|}
\toprule
Scenario & SAM slope & ESA slope & SAM / ESA ratio & SAM $R^2$ & ESA $R^2$ & Factor & Worse \\
\midrule
Disjoint Alphabet & 0.3588 & 0.3360 & 1.0677 & 0.9999 & 0.9999 & $1.07\times$ & SAM \\
\hline
Mutated Implant & 0.4070 & 0.2836 & 1.4349 & 1.0000 & 1.0000 & $1.43\times$ & SAM \\
\hline
Near Identical & 0.3404 & 0.3077 & 1.1061 & 0.9999 & 0.9999 & $1.11\times$ & SAM \\
\hline
Random Uniform & 0.3387 & 0.3556 & 0.9524 & 0.9998 & 0.9999 & $1.05\times$ & ESA \\
\hline
Repetitive With Noise & 0.4468 & 0.2824 & 1.5823 & 1.0000 & 1.0000 & $1.58\times$ & SAM \\
\bottomrule
\end{tabularx}
\end{benchmarktable}

